\section{Involution and Evolution}

The poet \textbf{John Keats} coined the term \textbf{Negative Capability}, that is, ``when a man is capable of being in uncertainties, mysteries, doubts, without any irritable reaching after fact and reason''.

To a higher type of mind, it may indicate the capacity to transcend easy answers and to consider things in their fullness. However, to a lesser mind it indicates confusion and inconsistency.

The modern mind is like the latter version as regards evolution. On the one hand, biological evolution proceeds without a plan, even without a cause, since random unpredictable variations is its motor. On the other hand, evolution is understood as the progression of the lower to the higher, although there is no scientific basis for such an assertion. In practice, moreover, that simply means that the current is ``higher'' than what came before. Hence, whatever modern men believe today is better than what they believed a thousand, a hundred, or even ten years ago.

Unfortunately, the corollary to that is that these current beliefs will in turn be outmoded. The ordinary person, and here I even include in this category those who are conventionally intelligent and educated, don't seem bothered by this. They assume they hold ``enlightened'' positions and are not concerned about the consequences. However, those who think more deeply see things much differently. For example, I saw a professor being interviewed about progressive politics. When presented with a list of names of those undeniably progressive of leftist, he refused to admit the obvious. After further interrogation, I realized that from the professor's point of view, none of those figures held opinions that a ``progressive'' will hold a generation from now. The conclusion I drew is that naïve progressive will be shocked by what they, or their children, will be expected to believe a 20 years from now.

I realize that intellectual moral arguments, such as natural law theory or Kant's categorical imperative, are well beyond the interests and capabilities of the vast majority, who just like to ``feel good'' about their moral positions. However, those who drive them and influence them are not unintelligent, since they already see the logical consequences, even if their followers do not. You can try this at home: try to extract the principle for some position, and draw out its applications to other domains. Ask a progressive about it. He will simply scoff at you; he will refute you with the claim that it is a stupid ``slippery slope'' argument. Nevertheless, who says A must say B, so eventually B will prevail.

\paragraph{Involution}


We recently provided an essay\footnote{\url{https://gornahoor.net/?p=6261}} by \textbf{Julius Evola} on his alternative theory to Darwinism, which was briefly reprised in the Revolt Against the Modern World. I don't know a traditional source for his precise presentation, although some concepts from Plotinus come through in it. It may also be a consequence of metaphysical positions. Obviously, past thinkers were hardly concerned with refuting Darwinism, and this would explain the lack of a precise argument. Therefore, Evola's ideas deserve a fair hearing as an alternative to Darwinism that doesn't depend on difficult to defend theories such as ``Creation Science''. It is clear that it is largely based on some concepts from Theosophy and Anthroposophy. \textbf{Edgar Dacqué} was a Theosophist. The ideas about animal ``group souls'' or of animals as ``cast offs'' from the human can be found in Rudolf Steiner's system. Perhaps they are found elsewhere in more traditional sources, but I am not aware of them.

However, we can and must point out a significant difference. It is not a question of spiritualizing matter, since that would require that matter be prior to the spirit as well as self-subsistent. No, prime matter, or unformed matter, does not exist. It is only when it is ``informed'' by spirit does it exist. Hence, no one can show you a piece of matter; there are stones, sticks, carbon atoms, etc., i.e., always something specific, but not matter in itself. This requires, too, the forms or ideas that inform matter. This reminds of a discussion\footnote{\url{https://gornahoor.net/?p=44}} I had several years ago with a proponent of Intelligent Design. Now ID wants to be a science, i.e., something than can be demonstrated empirically without recourse to any metaphysical systems. In that regard, it has not overturned Darwinism among scientists. However, as I was more interested in the metaphysics behind it, I asked the philosopher if ID commits one to a ``realist'' position, which of course I was referring to the Aristotelean-Thomist (A-T) system. The philosopher was taken aback by the question, but then decided it may be the case.

That is how to understand Evola's system. It presupposes that the idea of man is ontologically prior to any physical manifestation of man in time. Time is the projection of a transcendent idea into manifestation. Hence, what we experience as evolution in time is really the involution of the idea as it informs matter. Unlike the alternatives like Intelligent Design or Creation Science, the theory of involution is perfectly compatible with the current understanding of scientific data. Moreover, it also explains how ``evolution'', even though completely random, seems to arrive at man. Thus the Negative Capability is resolved in the higher synthesis: the twin understanding of Darwinian evolution as mentioned above is intelligently resolved.

However, involution is the enemy of egalitarianism. Some attempts will manifest closest to the ideal man, while others will be defective, incomplete, or abnormal. Further out, are the abortive lines that are the animals. Now individual animals are not intelligent (i.e., they do not have an intellectual soul), but there is a ``group soul'' for each animal species that guides its evolution. Hence, an individual animal is not the defective man, but rather it is the ``group soul''. When understood, that may overcome some objections that we've heard.

A weakness of the A-T system is that it only accounts for species, not individuals. Evola adds an interesting twist when he concludes that the individual is the form and the ``genus'' is the matter. This is worth further thought.

\paragraph{Genotypes and Involution}


The genotype is the genetic makeup of an organism. Let's assume, although of course we do not accept it as is, that the genotype accounts for all of an organism's characteristics. Moreover, the genotype relates an organism to other individuals of the same species. Mathematically, we can conceive of the set of all possible genotypes. Obviously, this will be an incredibly huge number, but that is of no significance. Simply put, this can be all possible DNA and RNA sequences, given some reasonable upper bound. Of this set, probably most are not viable, i.e., no life form can exist with those sequences. The remainder consists of possible life forms, i.e., they are possibilities of manifestation.

A species, then, can be understood as a particular subset of that set. For higher life forms, the prime operation is reproduction, or the combining of two different set members in a defined way to create a third. The first rule of a species subset is that it is closed under reproduction, i.e., reproduction will result in another member of the set, as long as there is no cross-breeding, hybridization, or random mutations. Hence, squirrels will always beget squirrels and Han Chinese will always beget Han Chinese. In this regard, species are quite stable.

This does not mean that the relative proportions of certain features cannot change over time as in the famous case of the black and white swallows; however, those variations are necessarily included in the species set and do not represent an innovation.

An innovation can only occur through cross-breeding, hybridization, or mutations from whatever causes. For example, the mating of two squirrels may result, if hit with a mutation that produces a viable genotype, in a new species, say a supersquirrel. Now, it is incorrect to say that the supersquirrel ``evolved'' from the squirrel, since the supersquirrel is not in the set of squirrel genotypes. Certainly the supersquirrels may displace the squirrels since they fight for the same food and do it better.

So now we can understand the process of involution better. The individual, depending on his spiritual characteristics, will attempt to incarnate as one of those species, i.e., as a particular genotype to which it has an affinity.

\paragraph{Retranslation of Passage from Revolt}


See page 180 of \textit{Revolt Against the Modern World}.

\begin{quotex}
Evolutionists believe they are ``positively'' keeping to the facts. They do not doubt that the facts in themselves are mute; the same facts, if interpreted in different ways, can lend support to the most varied hypotheses. It has happened, however, that someone, though having available of all the data adopted as proof of the theory of evolution, has shown that these data, in the final analysis, could support the opposite thesis, which in more than one respect corresponds to the traditional teaching. That is, the thesis that man is far from being a product of ``evolution'' of animal species, and many animal species must be considered as a lateral trunk in which the primordial impulse was aborted, having its direct and suitable manifestation only in the superior human races.
\end{quotex}


At this point, there is a reference to the works of Edgar Dacqué, \textit{The History of Involution} by E. Marconi, and \textit{The Transformist Illusion} by Dewar.


\flrightit{Posted on 2013-04-23 by Cologero}

\begin{center}* * *\end{center}

\begin{footnotesize}
\begin{sffamily}
\texttt{Matt on 2013-04-23 at 06:57 said:}

Cologero, this is a really good summary of the A-T influenced alternative to the neo-darwinian framework. And yes, so many self-described progressives still to this day will use that darwnian framework to support their ideas of social progress and the immanent eschaton. Without that failed recourse, one is only left with confused and incoherent answers from them as to just how this apparent social/historical progress occurs. No wonder post-modernism became so popular.

With that said, I just have a question in regards to the possibility of beings that have the material/biological shape of a human that are nevertheless animals -- lacking the capability of intelligence and reason, basically lacking the intellectual soul -- that one finds in a number of traditional doctrines (this was touched upon awhile back in the registered users forum). The possible result of cross-breeding (the answer usually seen in the traditional teachings)?

Sorry if my question comes across as vague. Maybe it should be best left for the forum.

\hfill

\texttt{Cologero on 2013-04-23 at 07:42 said:}

I believe that the anthropoids mentioned in the UR group article are such beings. As as you mention, it is part of traditional doctrine as a corollary to the idea of hierarchy. Mouravieff, in particular, proposes such a view.

However, that idea, which previously was only mentioned by ``fringe'' groups, is being discussed more, since it seems necessary to reconcile religious teachings with scientific data. For example, the Jewish physicist, Gerald Schroeder, proposes that idea (\textit{The Science of God}\footnote{\url{http://www.amazon.com/Science-God-Convergence-Scientific-Biblical/dp/1439129584/ref=la_B000APV1XA_1_1?ie=UTF8\&qid=1366715379\&sr=1-1}}. Closer to home, Edward Feser entertains the same idea (\textit{Monkey in your Soul}\footnote{\url{http://edwardfeser.blogspot.com/2011/09/monkey-in-your-soul.html}}. Now whether such anthropoids still exist is a touchy question.

\hfill

\texttt{Jason-Adam on 2013-04-23 at 13:46 said:}

Rather than attack science on their chosen ground with their assumptions, we need to go to the source and attack the scientific method itself. Man is a rational being so we can and must use reason and logic to find the truth. Thomism can be proved to be the correct philosophy and we need to stand up for it against modernism in all its forms.

\hfill

\texttt{Matt on 2013-04-23 at 20:47 said:}

I don't think this post (or the work of gornahoor, or tradition itself) is about attacking the work of scientists, or even the scientific/emprical method. It's to show that while a given theory can provide an empirically adequate description of observable phenomenon, it cannot provide an answer to certain important questions, such as question 2 in a previous post in this series of posts by Cologero. And it cannot provide these answers because the focus of the scientific method is not -- to use Aristotelian terminology -- on formal and final causes.

So what Evola was doing in those written pieces was to interpret the same data that biologists study from a metaphysical standpoint. I think Ananda Coomaramswamy gave a rather similar perspective on this subject -- thats if my memory serves me well (Cologero can correct me on this if it is not so).

\hfill

\texttt{Matt on 2013-04-23 at 20:51 said:}

Yes, that can be a touchy question... ..as seen by various ideologies in proclaiming entire historical ethnic groups as these non-human anthropoids.

\hfill

\texttt{Logres on 2013-04-24 at 13:53 said:}

Perhaps they are found elsewhere in more traditional sources, but I am not aware of them.''

I know Guenon says that Rosicrucian initiation is dead; although I am not sure if he or Evola were ever interested in Max Heindel -- in Heindel's Rosicrucian CosmoConception, he addresses group souls of more intelligent animals (such as horses, dogs, and cats), as well as the group ``waves'' of humanity; his cosmology was heavily influenced by evolutionary thinking, but the root of his thought is not evolutionary, in my opinion, because he assumes that Spirit dominates and informs matter, and in fact has an affinity more with Teilhard de Chardin's thought on the matter. At one stage in involution, he claims that man's spirit was ``outside'' man, and controlled the body, which is somewhat the case with the group souls that dominate animal species. He also addresses the ``race'' question, and asserts (in line I think with Traditional thought) that mere biological race is one of the ``traps'' for those spirits who seek to be free from matter, although he also admits that biological race is a definite vehicle : that is, different spirits have affinity with different races, to suit their purposes as they ``strive''. As an example, he cites the higher nervous development of the Aryan race wave, bought at the cost of some loss in physical hardihood in the physical body. The different waves are related to the planets, which govern each cycle.

\url{http://en.wikipedia.org/wiki/The_Rosicrucian_Cosmo-Conception}


\hfill

\texttt{Logres on 2013-04-24 at 14:01 said:}

JP Moreland wrote a fairly good work on the limits of scientific method:

``Christianity and the Nature of Science''. Nothing but group delusion can explain the irrational attachment to pseudo-science masquerading as Science which dominates modern thought. Men like EO Wilson predictably (as Cologero has written) attribute their own self-intuition of animal impulses to the entire species of humanity. What is true for them individually (but which they cannot, metaphysically, explain) is taken as true for everyone (eg., no one can have free will or act altruistically, or have metaphysical insight).

\url{http://en.wikipedia.org/wiki/E._O._Wilson}

Nevertheless, even men like this are ``censored'' by the liberal establishment for some implications in their worldview.

\end{sffamily}
\end{footnotesize}

